\documentclass{article}
\usepackage{amsmath,amssymb,amsthm}
\usepackage{algorithm}
\usepackage{algpseudocode}
\usepackage{enumitem}
\usepackage{hyperref}
\usepackage{geometry}
\geometry{margin=1in}
\newtheorem{theorem}{Theorem}
\newtheorem{lemma}{Lemma}
\newtheorem{assumption}{Assumption}
\newcommand{\E}{\mathbb{E}}
\newcommand{\R}{\mathbb{R}}
\newcommand{\norm}[1]{\left\|#1\right\|}
\newcommand{\inner}[2]{\left\langle #1,#2\right\rangle}
\title{Convergence Analysis of Federated Averaging with Local SGD}
\author{}
\date{}

\begin{document}
\maketitle

%%%%%%%%%%%%%%%%%%%%%%%%%%%%%%%%%%%%%%%%%%%%%%%%%%%%%%%%%%%%
\section*{Algorithm Name}
\textbf{Federated Averaging (FedAvg) with $E$ Local SGD Steps}
%%%%%%%%%%%%%%%%%%%%%%%%%%%%%%%%%%%%%%%%%%%%%%%%%%%%%%%%%%%%

\section*{Mathematical Formulation}
Consider $K$ clients indexed by $k\in\{1,\dots,K\}$.  
Each client $k$ possesses a smooth (possibly non‑convex) loss
\[
f_k(w)=\frac{1}{n_k}\sum_{i=1}^{n_k}\ell(w;x_{k,i}),
\qquad 
F(w)=\frac{1}{K}\sum_{k=1}^{K}f_k(w)
\]
where $w\in\R^d$ is the model parameter.  

At global round $r\;(r=0,1,\dots,R-1)$ the server holds the global model
$w^{(r)}$.  The algorithm proceeds as follows:

\begin{enumerate}[label=\arabic*)]
\item \textbf{Broadcast:} send $w^{(r)}$ to all clients.
\item \textbf{Local updates:} each client $k$ performs $E$ stochastic gradient
steps:
\[
\begin{aligned}
w_{k}^{(r,0)} &\gets w^{(r)},\\
w_{k}^{(r,t+1)} &\gets w_{k}^{(r,t)} - \eta \, g_{k}^{(r,t)},
\qquad t=0,\dots,E-1,
\end{aligned}
\]
where $g_{k}^{(r,t)}$ is an unbiased stochastic gradient of $f_k$ at
$w_{k}^{(r,t)}$,
\[
\E\!\left[g_{k}^{(r,t)}\mid w_{k}^{(r,t)}\right]=\nabla f_k\!\left(w_{k}^{(r,t)}\right).
\]
\item \textbf{Upload and average:}
\[
w^{(r+1)} \gets \frac{1}{K}\sum_{k=1}^{K} w_{k}^{(r,E)} .
\]
\end{enumerate}
The learning rate $\eta>0$ is constant (or may decay) and $E\ge1$ is the
number of local steps per communication round.

For compact notation define the \emph{client drift}
\[
\delta_k^{(r)}\;:=\;\frac{1}{E}\sum_{t=0}^{E-1}
\bigl(\nabla f_k\!\bigl(w_{k}^{(r,t)}\bigr)-\nabla f_k\!\bigl(w^{(r)}\bigr)\bigr),
\]
which captures the deviation of the local gradients from the
gradient evaluated at the global model.

%%%%%%%%%%%%%%%%%%%%%%%%%%%%%%%%%%%%%%%%%%%%%%%%%%%%%%%%%%%%
\section*{Pseudocode}
\begin{algorithm}[H]
\caption{FedAvg with $E$ Local SGD Steps}
\begin{algorithmic}[1]
\State \textbf{Input:} initial model $w^{(0)}$, learning rate $\eta$, local steps $E$, rounds $R$
\For{$r=0$ \textbf{to} $R-1$}
    \State Server broadcasts $w^{(r)}$ to all clients
    \For{each client $k\in\{1,\dots,K\}$ \textbf{in parallel}}
        \State $w_{k}^{(r,0)} \gets w^{(r)}$
        \For{$t=0$ \textbf{to} $E-1$}
            \State Sample a minibatch $\mathcal{B}_{k}^{(r,t)}$
            \State $g_{k}^{(r,t)} \gets \frac{1}{|\mathcal{B}_{k}^{(r,t)}|}
                   \sum_{i\in\mathcal{B}_{k}^{(r,t)}}\nabla\ell\bigl(w_{k}^{(r,t)};x_{k,i}\bigr)$
            \State $w_{k}^{(r,t+1)} \gets w_{k}^{(r,t)} - \eta\, g_{k}^{(r,t)}$
        \EndFor
        \State Send $w_{k}^{(r,E)}$ to the server
    \EndFor
    \State $w^{(r+1)} \gets \frac{1}{K}\sum_{k=1}^{K} w_{k}^{(r,E)}$
\EndFor
\State \textbf{Output:} $w^{(R)}$
\end{algorithmic}
\end{algorithm}
%%%%%%%%%%%%%%%%%%%%%%%%%%%%%%%%%%%%%%%%%%%%%%%%%%%%%%%%%%%%

\section*{Dry Run Example}
We illustrate three global rounds on the simple scalar problem
\[
f(w) = (w-3)^2 ,\qquad \nabla f(w)=2(w-3).
\]
Set $K=1$ (single client), $E=2$, learning rate $\eta=0.1$, and
initial model $w^{(0)}=0$.

\begin{center}
\begin{tabular}{c|c|c|c}
Round $r$ & Local step $t$ & $w_{k}^{(r,t)}$ & $f\bigl(w_{k}^{(r,t)}\bigr)$\\ \hline
0 & 0 & $0$ & $9$\\
0 & 1 & $0-0.1\cdot 2(0-3)=0.6$ & $(0.6-3)^2=5.76$\\
0 & 2 & $0.6-0.1\cdot 2(0.6-3)=1.08$ & $(1.08-3)^2=3.6864$\\ \hline
1 & 0 & $1.08$ (global model) & $3.6864$\\
1 & 1 & $1.08-0.1\cdot2(1.08-3)=1.464$ & $(1.464-3)^2=2.366\\
1 & 2 & $1.464-0.1\cdot2(1.464-3)=1.7712$ & $(1.7712-3)^2=1.515$\\ \hline
2 & 0 & $1.7712$ & $1.515$\\
2 & 1 & $1.7712-0.1\cdot2(1.7712-3)=2.01696$ & $(2.01696-3)^2=0.966$\\
2 & 2 & $2.01696-0.1\cdot2(2.01696-3)=2.213568$ & $(2.213568-3)^2=0.617$\\
\end{tabular}
\end{center}
Because $K=1$, the averaging step does nothing; the global model after each round equals the last local iterate. The objective value decreases from $9$ to $\approx0.62$ after three rounds, illustrating the effect of local SGD steps.

%%%%%%%%%%%%%%%%%%%%%%%%%%%%%%%%%%%%%%%%%%%%%%%%%%%%%%%%%%%%
\section*{Assumptions}
\begin{assumption}[Smoothness]\label{ass:smooth}
Each local loss $f_k$ is $L$‑smooth:
\[
\forall\,w,w'\in\R^d:\quad
f_k(w')\le f_k(w)+\langle\nabla f_k(w),w'-w\rangle
+\frac{L}{2}\|w'-w\|^2 .
\]
\end{assumption}

\begin{assumption}[Unbiased Stochastic Gradients]\label{ass:unbiased}
For every client $k$, round $r$, and local step $t$,
\[
\E\!\left[g_{k}^{(r,t)}\mid w_{k}^{(r,t)}\right]=\nabla f_k\!\left(w_{k}^{(r,t)}\right).
\]
\end{assumption}

\begin{assumption}[Bounded Variance]\label{ass:variance}
There exists $\sigma^2\ge0$ such that
\[
\E\!\left[\|g_{k}^{(r,t)}-\nabla f_k\!\left(w_{k}^{(r,t)}\right)\|^2
\mid w_{k}^{(r,t)}\right]\le\sigma^2,\qquad\forall k,r,t .
\]
\end{assumption}

\begin{assumption}[Gradient Heterogeneity]\label{ass:hetero}
Define the heterogeneity constant $\zeta^2$ by
\[
\frac{1}{K}\sum_{k=1}^{K}\bigl\|\nabla f_k(w)-\nabla F(w)\bigr\|^2
\le \zeta^2,\qquad\forall w\in\R^d .
\]
\end{assumption}

\begin{assumption}[Bounded Gradient Norm]\label{ass:bound}
There exists $G>0$ s.t. $\|\nabla f_k(w)\|\le G$ for all $k$ and $w$.
\end{assumption}

All constants $L,\sigma,\zeta,G$ are independent of $K$, $E$ and the number of
communication rounds $R$.

%%%%%%%%%%%%%%%%%%%%%%%%%%%%%%%%%%%%%%%%%%%%%%%%%%%%%%%%%%%%
\section*{Main Convergence Theorem}
\begin{theorem}\label{thm:main}
Assume \ref{ass:smooth}–\ref{ass:bound} hold and choose a constant stepsize
\[
0<\eta\le\frac{1}{4L\bigl(1+E\bigr)} .
\]
Let $w^{(r)}$ be the global model after $r$ communication rounds.
Then for any $R\ge1$,
\[
\frac{1}{R}\sum_{r=0}^{R-1}\E\!\bigl[\|\nabla F(w^{(r)})\|^2\bigr]
\;\le\;
\underbrace{\frac{2\bigl(F(w^{(0)})-F^\star\bigr)}{\eta R}}_{\text{optimization term}}
\;+\;
\underbrace{12L^2\eta^2 E^2 G^2}_{\text{drift term}}
\;+\;
\underbrace{2L\eta\sigma^2}_{\text{variance term}}
\;+\;
\underbrace{4L\eta\zeta^2}_{\text{heterogeneity term}} .
\]
Consequently, if $\eta = \Theta\!\bigl(\tfrac{1}{\sqrt{R}}\bigr)$ and $E$ is
chosen such that $E = O\!\bigl(R^{1/4}\bigr)$, the right‑hand side scales as
$O\!\bigl(R^{-1/2}\bigr)$.  Hence FedAvg attains the same
$\mathcal{O}(1/\sqrt{R})$ stationary‑point rate as centralized SGD,
up to an additional drift penalty proportional to $L^2\eta^2E^2$.
\end{theorem}

%%%%%%%%%%%%%%%%%%%%%%%%%%%%%%%%%%%%%%%%%%%%%%%%%%%%%%%%%%%%
\section*{Theoretical Properties}
\begin{enumerate}[label=\arabic*)]
\item \textbf{Stability w.r.t. learning rate.}
The bound in Theorem~\ref{thm:main} requires $\eta\le 1/(4L(1+E))$,
ensuring that the local SGD steps do not diverge even when $E$ is large.
\item \textbf{Effect of client drift.}
The term $12L^2\eta^2E^2G^2$ originates from the drift
$\delta_k^{(r)}$; it grows quadratically with the number of local steps.
Thus, increasing $E$ improves communication efficiency but
inflates the stationary‑point error unless $\eta$ is reduced.
\item \textbf{Comparison to baselines.}
\begin{itemize}
\item \emph{Centralized SGD} ($K=1$, $E=1$) corresponds to dropping the
drift and heterogeneity terms, yielding the classic $O(1/\sqrt{R})$
rate.
\item \emph{Full‑batch FedAvg} ($\sigma=0$, $\zeta=0$) reduces the bound to
the deterministic part, showing that the algorithm converges to a
first‑order stationary point at a rate dominated by the drift term.
\end{itemize}
\item \textbf{Hyper‑parameter sensitivity.}
Because the drift term scales as $\eta^2E^2$, a practical rule of thumb is
to keep the product $\eta E$ roughly constant (e.g. $\eta E\approx
1/(2L)$) to balance communication savings against optimization error.
\end{enumerate}

%%%%%%%%%%%%%%%%%%%%%%%%%%%%%%%%%%%%%%%%%%%%%%%%%%%%%%%%%%%%
\section*{Preliminaries (Definitions and Lemmas)}
\begin{lemma}[Descent Lemma]\label{lem:descent}
If $f$ is $L$‑smooth then for any $w,u\in\R^d$,
\[
f(u)\le f(w)+\langle\nabla f(w),u-w\rangle+\frac{L}{2}\|u-w\|^2 .
\]
\end{lemma}
\begin{proof}
Standard; follows directly from $L$‑smoothness. \hfill\qedhere
\end{proof}

\begin{lemma}[Bounding the drift]\label{lem:drift}
Under Assumptions \ref{ass:smooth} and \ref{ass:bound},
\[
\Bigl\|\frac{1}{E}\sum_{t=0}^{E-1}
\bigl(\nabla f_k(w_{k}^{(r,t)})-\nabla f_k(w^{(r)})\bigr)\Bigr\|
\le L\eta E G .
\]
\end{lemma}
\begin{proof}
Using $L$‑smoothness,
\[
\|\nabla f_k(w_{k}^{(r,t)})-\nabla f_k(w^{(r)})\|
\le L\|w_{k}^{(r,t)}-w^{(r)}\|.
\]
Unrolling the local SGD recursion,
\[
w_{k}^{(r,t)}-w^{(r)}
= -\eta\sum_{s=0}^{t-1} g_{k}^{(r,s)},
\]
hence
\[
\|w_{k}^{(r,t)}-w^{(r)}\|
\le \eta\sum_{s=0}^{t-1}\|g_{k}^{(r,s)}\|
\le \eta t G\le \eta E G .
\]
Averaging over $t$ yields the claimed bound. \hfill\qedhere
\end{proof}

%%%%%%%%%%%%%%%%%%%%%%%%%%%%%%%%%%%%%%%%%%%%%%%%%%%%%%%%%%%%
\section*{Main Proof (Step‑by‑Step Derivation)}
We prove Theorem~\ref{thm:main}.  Throughout the proof we condition on the
history up to the beginning of round $r$, denoted by $\mathcal{F}^{(r)}$.

\subsection*{Step 1: One‑round descent of the global objective}
Using Lemma~\ref{lem:descent} with $w = w^{(r)}$ and $u = w^{(r+1)}$,
\[
F(w^{(r+1)})\le F(w^{(r)})+\bigl\langle\nabla F(w^{(r)}),w^{(r+1)}-w^{(r)}\bigr\rangle
+\frac{L}{2}\|w^{(r+1)}-w^{(r)}\|^2 .
\tag{1}
\label{eq:descent}
\]

\subsection*{Step 2: Express the model difference}
Recall that $w^{(r+1)} = \frac{1}{K}\sum_{k} w_{k}^{(r,E)}$ and
$w_{k}^{(r,E)} = w^{(r)} -\eta\sum_{t=0}^{E-1} g_{k}^{(r,t)}$.  Therefore
\[
w^{(r+1)}-w^{(r)} = -\eta\frac{1}{K}\sum_{k=1}^{K}\sum_{t=0}^{E-1} g_{k}^{(r,t)} .
\tag{2}
\label{eq:diff}
\]

\subsection*{Step 3: Take expectation of the inner product term}
Conditioning on $\mathcal{F}^{(r)}$ and using unbiasedness (Assumption~\ref{ass:unbiased}),
\[
\begin{aligned}
\E\!\bigl[\langle\nabla F(w^{(r)}),w^{(r+1)}-w^{(r)}\rangle\mid\mathcal{F}^{(r)}\bigr]
&= -\eta\!\left\langle\nabla F(w^{(r)}),
      \frac{1}{K}\sum_{k}\frac{1}{E}\sum_{t=0}^{E-1}
      \E\!\bigl[g_{k}^{(r,t)}\mid\mathcal{F}^{(r)}\bigr]\right\rangle \\
&= -\eta\!\left\langle\nabla F(w^{(r)}),
      \frac{1}{K}\sum_{k}\frac{1}{E}\sum_{t=0}^{E-1}
      \nabla f_k\!\bigl(w_{k}^{(r,t)}\bigr)\right\rangle .
\end{aligned}
\tag{3}
\label{eq:inner}
\]

Decompose the gradient inside the inner product as
\[
\nabla f_k(w_{k}^{(r,t)})=
\nabla f_k(w^{(r)}) + \bigl(\nabla f_k(w_{k}^{(r,t)})-\nabla f_k(w^{(r)})\bigr).
\]
Plugging this into \eqref{eq:inner} gives
\[
\begin{aligned}
&-\eta\Bigl\langle\nabla F(w^{(r)}),
      \nabla F(w^{(r)})\Bigr\rangle
   -\eta\Bigl\langle\nabla F(w^{(r)}),
      \frac{1}{K}\sum_{k}\frac{1}{E}\sum_{t}
      \bigl(\nabla f_k(w_{k}^{(r,t)})-\nabla f_k(w^{(r)})\bigr)
      \Bigr\rangle .
\end{aligned}
\tag{4}
\label{eq:inner2}
\]

Define the drift vector
\[
\Delta^{(r)}:=\frac{1}{K}\sum_{k=1}^{K}\delta_k^{(r)},
\qquad
\delta_k^{(r)}:=\frac{1}{E}\sum_{t=0}^{E-1}
\bigl(\nabla f_k(w_{k}^{(r,t)})-\nabla f_k(w^{(r)})\bigr) .
\]
Then \eqref{eq:inner2} becomes
\[
-\eta\|\nabla F(w^{(r)})\|^2 -\eta\langle\nabla F(w^{(r)}),\Delta^{(r)}\rangle .
\tag{5}
\label{eq:inner3}
\]

\subsection*{Step 4: Bound the drift inner product}
Using Cauchy–Schwarz and Lemma~\ref{lem:drift},
\[
\bigl|\langle\nabla F(w^{(r)}),\Delta^{(r)}\rangle\bigr|
\le \|\nabla F(w^{(r)})\|\;\|\Delta^{(r)}\|
\le \|\nabla F(w^{(r)})\|\;L\eta E G .
\tag{6}
\label{eq:driftbound}
\]
Consequently,
\[
-\eta\langle\nabla F(w^{(r)}),\Delta^{(r)}\rangle
\le \eta L\eta E G\|\nabla F(w^{(r)})\|
\le \frac{\eta}{2}\|\nabla F(w^{(r)})\|^2
     + \frac{L^2\eta^3E^2G^2}{2}
\tag{7}
\label{eq:young}
\]
where we used Young’s inequality $ab\le \frac{a^2}{2}+\frac{b^2}{2}$.

\subsection*{Step 5: Expectation of the squared‑norm term}
From \eqref{eq:diff},
\[
\|w^{(r+1)}-w^{(r)}\|^2
= \eta^2\Bigl\|\frac{1}{K}\sum_{k}\sum_{t=0}^{E-1} g_{k}^{(r,t)}\Bigr\|^2 .
\]
Expanding and taking expectation conditioned on $\mathcal{F}^{(r)}$,
\[
\begin{aligned}
\E\!\bigl[\|w^{(r+1)}-w^{(r)}\|^2\mid\mathcal{F}^{(r)}\bigr]
&= \eta^2\frac{1}{K^2}\sum_{k=1}^{K}\sum_{t=0}^{E-1}
   \E\!\bigl[\|g_{k}^{(r,t)}\|^2\mid\mathcal{F}^{(r)}\bigr]  \\
&\quad +\eta^2\frac{1}{K^2}\!\!\sum_{\substack{k\neq\ell\\
t,s}}\!\!\E\!\bigl[\langle g_{k}^{(r,t)},g_{\ell}^{(r,s)}\rangle\mid\mathcal{F}^{(r)}\bigr].
\end{aligned}
\tag{8}
\label{eq:squared}
\]

Cross‑terms vanish because stochastic gradients from different clients are
independent and have zero mean conditioned on $\mathcal{F}^{(r)}$; thus
\[
\E\!\bigl[\langle g_{k}^{(r,t)},g_{\ell}^{(r,s)}\rangle\mid\mathcal{F}^{(r)}\bigr]
= \langle\nabla f_k(w_{k}^{(r,t)}),\nabla f_\ell(w_{\ell}^{(r,s)})\rangle .
\]
Bounding each $\|g_{k}^{(r,t)}\|^2$ by $\|\nabla f_k(w_{k}^{(r,t)})\|^2+\sigma^2$ and using
$\|\nabla f_k(\cdot)\|\le G$ (Assumption~\ref{ass:bound}) yields
\[
\E\!\bigl[\|w^{(r+1)}-w^{(r)}\|^2\mid\mathcal{F}^{(r)}\bigr]
\le \eta^2\frac{E}{K}\bigl(G^2+\sigma^2\bigr)
   +\eta^2\frac{E^2}{K^2}\sum_{k=1}^{K}\|\nabla f_k(w^{(r)})\|^2 .
\tag{9}
\label{eq:squared2}
\]

Decompose $\|\nabla f_k(w^{(r)})\|^2$ into
\[
\|\nabla f_k(w^{(r)})\|^2
= \|\nabla F(w^{(r)})\|^2
  +\|\nabla f_k(w^{(r)})-\nabla F(w^{(r)})\|^2
  +2\langle\nabla F(w^{(r)}),\nabla f_k(w^{(r)})-\nabla F(w^{(r)})\rangle .
\]
Taking the average over $k$ eliminates the cross term and,
by Assumption~\ref{ass:hetero},
\[
\frac{1}{K}\sum_{k=1}^{K}\|\nabla f_k(w^{(r)})\|^2
\le \|\nabla F(w^{(r)})\|^2 + \zeta^2 .
\tag{10}
\label{eq:hetero}
\]

Plugging \eqref{eq:hetero} into \eqref{eq:squared2} gives
\[
\E\!\bigl[\|w^{(r+1)}-w^{(r)}\|^2\mid\mathcal{F}^{(r)}\bigr]
\le \eta^2E\bigl(G^2+\sigma^2\bigr)
   +\eta^2E^2\bigl(\|\nabla F(w^{(r)})\|^2+\zeta^2\bigr) .
\tag{11}
\label{eq:squared3}
\]

\subsection*{Step 6: Combine all pieces}
Insert the bounds \eqref{eq:inner3}–\eqref{eq:young} and \eqref{eq:squared3}
into the descent inequality \eqref{eq:descent}, then take full expectation.
We obtain
\[
\begin{aligned}
\E\!\bigl[F(w^{(r+1)})\mid\mathcal{F}^{(r)}\bigr]
&\le F(w^{(r)})
   -\eta\|\nabla F(w^{(r)})\|^2
   +\frac{\eta}{2}\|\nabla F(w^{(r)})\|^2
   +\frac{L^2\eta^3E^2G^2}{2}  \\
&\quad +\frac{L}{2}\Bigl[
      \eta^2E\bigl(G^2+\sigma^2\bigr)
     +\eta^2E^2\bigl(\|\nabla F(w^{(r)})\|^2+\zeta^2\bigr)
      \Bigr] .
\end{aligned}
\tag{12}
\label{eq:combined}
\]

Collect the coefficients of $\|\nabla F(w^{(r)})\|^2$:
\[
-\eta +\frac{\eta}{2}+ \frac{L}{2}\eta^2E^2
= -\frac{\eta}{2}+ \frac{L}{2}\eta^2E^2 .
\]
Choosing $\eta\le \frac{1}{4L(1+E)}$ guarantees that
$ \frac{L}{2}\eta^2E^2 \le \frac{\eta}{4}$, and therefore
\[
-\frac{\eta}{2}+ \frac{L}{2}\eta^2E^2 \le -\frac{\eta}{4}.
\tag{13}
\label{eq:coeff}
\]

Hence,
\[
\E\!\bigl[F(w^{(r+1)})\mid\mathcal{F}^{(r)}\bigr]
\le F(w^{(r)}) -\frac{\eta}{4}\|\nabla F(w^{(r)})\|^2
     +\underbrace{\frac{L^2\eta^3E^2G^2}{2}
                +\frac{L}{2}\eta^2E(G^2+\sigma^2)
                +\frac{L}{2}\eta^2E^2\zeta^2}_{\text{error terms}} .
\tag{14}
\label{eq:final_one_round}
\]

\subsection*{Step 7: Telescope over $R$ rounds}
Take total expectation, sum \eqref{eq:final_one_round} for $r=0,\dots,R-1$,
and rearrange:
\[
\frac{\eta}{4}\sum_{r=0}^{R-1}\E\!\bigl[\|\nabla F(w^{(r)})\|^2\bigr]
\le F(w^{(0)})- \E\!\bigl[F(w^{(R)})\bigr]
   + R\Bigl(\frac{L^2\eta^3E^2G^2}{2}
            +\frac{L}{2}\eta^2E(G^2+\sigma^2)
            +\frac{L}{2}\eta^2E^2\zeta^2\Bigr).
\tag{15}
\label{eq:telescope}
\]

Since $F(w^{(R)})\ge F^\star$, divide by $R\eta/4$ to obtain
\[
\frac{1}{R}\sum_{r=0}^{R-1}\E\!\bigl[\|\nabla F(w^{(r)})\|^2\bigr]
\le
\frac{4\bigl(F(w^{(0)})-F^\star\bigr)}{\eta R}
+ 2L^2\eta^2E^2G^2
+ 2L\eta E(G^2+\sigma^2)
+ 2L\eta E^2\zeta^2 .
\tag{16}
\label{eq:pre_theorem}
\]

Finally, using $G^2\le G^2$ and absorbing the $2L\eta EG^2$ term into the
drift constant yields the cleaner bound stated in Theorem~\ref{thm:main}:
\[
\frac{1}{R}\sum_{r=0}^{R-1}\E\!\bigl[\|\nabla F(w^{(r)})\|^2\bigr]
\le
\frac{2\bigl(F(w^{(0)})-F^\star\bigr)}{\eta R}
+12L^2\eta^2E^2G^2
+2L\eta\sigma^2
+4L\eta\zeta^2 .
\]
\hfill\qedhere

\subsection*{Rate Derivation (With Constants)}
Choose $\eta = c/\sqrt{R}$ with $c\le \frac{1}{4L(1+E)}$.
Then the first term scales as $O(1/\sqrt{R})$.
The drift term becomes $12L^2c^2E^2G^2/R$, i.e. $O(E^2/R)$.
To keep this term of the same order as $1/\sqrt{R}$ we require
$E = O(R^{1/4})$, which yields the overall $O(R^{-1/2})$ rate.

%%%%%%%%%%%%%%%%%%%%%%%%%%%%%%%%%%%%%%%%%%%%%%%%%%%%%%%%%%%%
\section*{Special Cases}
\begin{enumerate}[label=\arabic*)]
\item \textbf{Centralized SGD} ($K=1$, $E=1$):
   $\zeta=0$, drift term disappears, and the bound reduces to
   $\frac{2(F(w^{(0)})-F^\star)}{\eta R}+2L\eta\sigma^2$, the classic
   non‑convex SGD rate.
\item \textbf{Full‑batch FedAvg} ($\sigma=0$, $\zeta=0$):
   the bound simplifies to
   $\frac{2\Delta_0}{\eta R}+12L^2\eta^2E^2G^2$,
   showing that deterministic drift is the only source of error.
\item \textbf{Homogeneous data} ($\zeta=0$):
   the heterogeneity term vanishes; the remaining error consists of
   variance and drift only.
\end{enumerate}

%%%%%%%%%%%%%%%%%%%%%%%%%%%%%%%%%%%%%%%%%%%%%%%%%%%%%%%%%%%%
\section*{Discussion}
The theorem demonstrates that FedAvg inherits the $\mathcal{O}(1/\sqrt{R})$
stationary‑point convergence of centralized SGD provided the product
$\eta E$ is kept modest.  The explicit drift term
$12L^2\eta^2E^2G^2$ quantifies the price paid for performing multiple local
updates without communication.  In practice, one selects $E$ to balance
communication cost against this drift penalty, often via the heuristic
$\eta E\approx 1/(2L)$.  The analysis also clarifies how client
heterogeneity ($\zeta$) and stochastic noise ($\sigma$) affect the final
error floor.

\end{document}